% \looseness ASKS FOR A PARAGRAPH SO MANY LINES LONGER THAN THE BEST.
%
% The reference takes the way of breaking whose line count comes closest
% to best+\looseness, and among equals the cheapest. A paragraph that
% breaks into three lines:
%
%   \looseness=0    3 lines
%   \looseness=1    4
%   \looseness=2    5
%   \looseness=9   12
%   \looseness=-1   3      it cannot be made shorter, so it is not
%
% hstex ignored the parameter and always took the cheapest, so every one
% of those was three.
%
% IT ALSO CHANGES WHICH BREAKS ARE ALIKE. With any \looseness at all,
% every line number is its own class, whatever \parshape and \hangindent
% say -- see
% tests/trip/probes/every-break-is-weighed-against-its-own-line.tex --
% because the line count is now part of the answer and cannot be
% collapsed away.
%
% AND WHAT CANNOT BE HAD IS TRIED FOR AGAIN. If the asked-for looseness
% is not reached, the reference goes round another pass with more
% tolerance rather than settling; only the final pass takes what it can
% get. trip line 208 sets \looseness 5 over a paragraph that can only
% just manage it.
%
% \looseness belongs to one paragraph: \par gives it back to nought.
\tracingonline=1 \hsize=100pt \parindent=0pt \tolerance=10000 \hbadness=10000
\def\body{aa bb cc dd ee ff gg hh ii jj kk ll mm nn oo pp qq rr ss tt uu vv}
\message{[A looseness=0]}
\setbox0=\vbox{\noindent\body\par\message{[lines=\the\prevgraf]}}
\message{[B looseness=1]}
\setbox1=\vbox{\looseness=1 \noindent\body\par\message{[lines=\the\prevgraf]}}
\message{[C looseness=2]}
\setbox2=\vbox{\looseness=2 \noindent\body\par\message{[lines=\the\prevgraf]}}
\message{[D looseness=-1]}
\setbox3=\vbox{\looseness=-1 \noindent\body\par\message{[lines=\the\prevgraf]}}
\message{[E looseness=9]}
\setbox4=\vbox{\looseness=9 \noindent\body\par\message{[lines=\the\prevgraf]}}
\message{[F looseness is given back]}
\setbox5=\vbox{\looseness=1 \noindent\body\par\message{[after=\the\looseness]}}
\message{[done]}
\end
